\documentclass[10pt]{article}
\usepackage[utf8]{inputenc}
\usepackage[T1]{fontenc}
\usepackage{amsmath}
\usepackage{amsfonts}
\usepackage{amssymb}
\usepackage[version=4]{mhchem}
\usepackage{stmaryrd}
\usepackage{bbold}
\usepackage{hyperref}
\hypersetup{colorlinks=true, linkcolor=blue, filecolor=magenta, urlcolor=cyan,}
\urlstyle{same}

\title{On a construction of quadratic APN functions }

\author{Lilya Budaghyan, Claude Carlet, and Gregor Leander}
\date{}


%New command to display footnote whose markers will always be hidden
\let\svthefootnote\thefootnote
\newcommand\blfootnotetext[1]{%
  \let\thefootnote\relax\footnote{#1}%
  \addtocounter{footnote}{-1}%
  \let\thefootnote\svthefootnote%
}

%Overriding the \footnotetext command to hide the marker if its value is `0`
\let\svfootnotetext\footnotetext
\renewcommand\footnotetext[2][?]{%
  \if\relax#1\relax%
    \ifnum\value{footnote}=0\blfootnotetext{#2}\else\svfootnotetext{#2}\fi%
  \else%
    \if?#1\ifnum\value{footnote}=0\blfootnotetext{#2}\else\svfootnotetext{#2}\fi%
    \else\svfootnotetext[#1]{#2}\fi%
  \fi
}

\begin{document}
\maketitle


\begin{abstract}
In a recent paper, the authors introduced a method for constructing new quadratic APN functions from known ones. Applying this method, they obtained the function $x^{3}+\operatorname{tr}_{n}\left(x^{9}\right)$ which is APN over $\mathbb{F}_{2^{n}}$ for any positive integer $n$. The present paper is a continuation of this work. We give sufficient conditions on linear functions $L_{1}$ and $L_{2}$ from $\mathbb{F}_{2^{n}}$ to itself such that the function $L_{1}\left(x^{3}\right)+L_{2}\left(x^{9}\right)$ is APN over $\mathbb{F}_{2^{n}}$. We show that this can lead to many new cases of APN functions. In particular, we get two families of APN functions $x^{3}+a^{-1} \operatorname{tr}_{n}^{3}\left(a^{3} x^{9}+a^{6} x^{18}\right)$ and $x^{3}+a^{-1} \operatorname{tr}_{n}^{3}\left(a^{6} x^{18}+a^{12} x^{36}\right)$ over $\mathbb{F}_{2^{n}}$ for any $n$ divisible by 3 and $a \in \mathbb{F}_{2^{n}}^{*}$. We prove that for $n=9$, these families are pairwise different and differ from all previously known families of APN functions, up to the most general equivalence notion, the CCZ-equivalence. We also investigate further sufficient conditions under which the conditions on the linear functions $L_{1}$ and $L_{2}$ are satisfied.
\end{abstract}

Index Terms-Almost bent, Almost perfect nonlinear, CCZequivalence, Nonlinearity, S-box, Vectorial Boolean function

\section*{I. Introduction}
For any positive integer $n$, a function $F$ from the field $\mathbb{F}_{2^{n}}$ to itself is called almost perfect nonlinear (APN) if for every $a \neq 0$ and every $b$ in $\mathbb{F}_{2^{n}}$, the equation $F(x)+F(x+a)= b$ admits at most 2 solutions [9]. When used as S -boxes in block ciphers, APN functions oppose an optimum resistance to differential cryptanalysis (see [1]). The notion of APN function is closely connected to the notion of almost bent (AB) function [6], which can be described in terms of the Walsh transform of functions. Let $\operatorname{tr}_{n}(x)$ denote the trace function from $\mathbb{F}_{2^{n}}$ into $\mathbb{F}_{2}$, that is, $\operatorname{tr}_{n}(x)=x+x^{2}+x^{4}+\ldots+x^{2^{n-1}}$ - more generally, for any positive divisor $m$ of $n$, we shall also use the notation $\operatorname{tr}_{n}^{m}(x)=x+x^{2^{m}}+x^{2^{2 m}}+\ldots+x^{2^{(n / m-1) m}}-$ then the Walsh transform of a given function $F: \mathbb{F}_{2^{n}} \rightarrow \mathbb{F}_{2^{n}}$ is the integer valued function

$$
\lambda_{F}(a, b)=\sum_{x \in \mathbb{F}_{2^{n}}}(-1)^{\operatorname{tr}_{n}(b F(x)+a x)}, \quad a, b \in \mathbb{F}_{2^{n}} .
$$

Note that the Walsh transform can also be defined in the framework of functions from the vector space $\mathbb{F}_{2}^{n}$ to itself (or to another vector space) but we shall always endow here the vector space $\mathbb{F}_{2}^{n}$ with the structure of the field $\mathbb{F}_{2^{n}}$.\\
The set $\Lambda_{F}=\left\{\lambda_{F}(a, b): a, b \in \mathbb{F}_{2^{n}}, b \neq 0\right\}$ is called the Walsh spectrum of $F$ and the set $\left\{\left|\lambda_{F}(a, b)\right|: a, b \in \mathbb{F}_{2^{n}}, b \neq\right. 0\}$ the extended Walsh spectrum of $F$. If the Walsh spectrum of $F$ equals $\left\{0, \pm 2^{\frac{n+1}{2}}\right\}$ then the function $F$ is called AB [6]. AB functions exist for $n$ odd only and oppose an optimum

\footnotetext{Department of Informatics, University of Bergen, PB 7803, 5020 Bergen, NORWAY; e-mail: \href{mailto:Lilya.Budaghyan@ii.uib.no}{Lilya.Budaghyan@ii.uib.no}

Universities of Paris 8 and Paris 13; CNRS, UMR 7539 LAGA; Address: University of Paris 8, Department of Mathematics, 2 rue de la liberté, 93526 Saint-Denis cedex 02, FRANCE; e-mail: \href{mailto:claude.carlet@inria.fr}{claude.carlet@inria.fr}

Department of Mathematics, Technical University of Denmark, DENMARK; e-mail: \href{mailto:gregor.leander@rub.de}{gregor.leander@rub.de}
}
resistance to linear cryptanalysis (see [8]). Besides, every AB function is APN [6], and in the $n$ odd case, any quadratic function is APN if and only if it is AB [5]. A comprehensive survey on APN and AB functions can be found in [4].

The APN and AB properties are preserved by some transformations of functions. If $F$ is an APN (respectively AB ) function, $A_{1}, A_{2}$ are affine permutations and $A$ is an affine function, then the function $F^{\prime}=A_{1} \circ F \circ A_{2}+A$ is also APN [9] (respectively AB ). The functions $F$ and $F^{\prime}$ are then called extended affine equivalent (EA-equivalent). In [5], Carlet, Charpin and Zinoviev introduced an equivalence relation of functions, more recently called CCZ-equivalence, which preserves APN and AB properties: two mappings $F$ and $F^{\prime}$ from $\mathbb{F}_{2^{n}}$ to itself are called CCZ-equivalent if by some affine permutation $\mathcal{L}$ of $\mathbb{F}_{2^{n}}^{2}$ the image of the graph of $F$ is the graph of $F^{\prime}$, that is, $\mathcal{L}\left(G_{F}\right)=G_{F^{\prime}}$ where $G_{F}= \left\{(x, F(x)) \mid x \in \mathbb{F}_{2^{n}}\right\}$ and $G_{F^{\prime}}=\left\{\left(x, F^{\prime}(x)\right) \mid x \in \mathbb{F}_{2^{n}}\right\}$. Two CCZ-equivalent functions have the same extended Walsh spectrum and are both $\mathrm{APN} / \mathrm{AB}$ or not. EA-equivalence is a particular case of CCZ-equivalence and any permutation is CCZ-equivalent to its inverse. In [2], it is proven that CCZequivalence is more general than the conjunction of the EAequivalence of functions and the inversion of permutations.

In the recent paper [3], it is proven that the function $x^{3}+ \operatorname{tr}_{n}\left(x^{9}\right)$ is APN over $\mathbb{F}_{2^{n}}$ for any $n$, and that for almost all $n \geq$ 7 it is CCZ-inequivalent to the Gold functions (and therefore, being quadratic, EA-inequivalent to power functions), to the inverse and Dobbertin functions. And (being quadratic) this function is AB for every odd $n$. The present paper generalizes the construction

$$
x^{3}+\operatorname{tr}_{n}\left(x^{9}\right)
$$

to the form

$$
F(x)=L_{1}\left(x^{3}\right)+L_{2}\left(x^{9}\right)
$$

where $L_{1}$ and $L_{2}$ are linear functions from $\mathbb{F}_{2^{n}}$ to itself, and it studies conditions on $L_{1}$ and $L_{2}$ such that $F$ is APN. In particular, we prove that, if $n$ is even and the function $L_{1}(x)+ L_{2}\left(x^{3}\right)$ is a permutation of $\mathbb{F}_{2^{n}}$, then $F$ is APN. We have an example of the fact that even slightly changing functions $L_{1}$ and $L_{2}$ can lead to new cases of APN functions: for $n=6,8$, and $a$ a primitive element of $\mathbb{F}_{2^{n}}^{*}$, the APN functions $x^{3}+ \operatorname{tr}_{n}\left(x^{9}\right)$ and $a x^{3}+\operatorname{tr}_{n}\left(a^{3} x^{9}\right)$ are CCZ-inequivalent. Besides, we construct a few families of APN functions: let $n$ be a positive integer and $a \in \mathbb{F}_{2^{n}}^{*}$; the function

$$
F_{1}(x)=x^{3}+a^{-1} \operatorname{tr}_{n}\left(a^{3} x^{9}\right)
$$

and for $n$ divisible by 3 the functions

$$
\begin{aligned}
& F_{2}(x)=x^{3}+a^{-1} \operatorname{tr}_{n}^{3}\left(a^{6} x^{18}+a^{12} x^{36}\right), \\
& F_{3}(x)=x^{3}+a^{-1} \operatorname{tr}_{n}^{3}\left(a^{3} x^{9}+a^{6} x^{18}\right)
\end{aligned}
$$

are APN over $\mathbb{F}_{2^{n}}$. As mentioned above, the function $F_{1}$ is CCZ-inequivalent to $x^{3}+\operatorname{tr}_{n}\left(x^{9}\right)$ when $n=6,8$ and $a$ is a primitive element of $\mathbb{F}_{2^{6}}^{*}$, and, therefore, we can see that $F_{1}$ defines a new class of APN functions which includes the class of functions $x^{3}+\operatorname{tr}_{n}\left(x^{9}\right)$. Besides, we show that for $n=9$, the functions $F_{1}, F_{2}$ and $F_{3}$ are pairwise CCZ-inequivalent, and the functions $F_{2}$ and $F_{3}$ are CCZ-inequivalent to any function from previously known families of APN functions. Hence, $F_{1}$ generalizes the family of functions $x^{3}+\operatorname{tr}_{n}\left(x^{9}\right)$ while $F_{2}$ and $F_{3}$ define two new families of APN functions.

\section*{II. A construction of APN functions}
\section*{A. The case $n$ even}
There is a sufficient condition for a function in even number of variables of the form $x^{3}+L\left(x^{9}\right)$ (where $L$ is linear) to be APN, which explains many cases of APN functions of this form.

Proposition 1: Let $n$ be an even positive integer and $L$ a linear function from $\mathbb{F}_{2^{n}}$ to itself. If the function

$$
F(x)=x+L\left(x^{3}\right)
$$

is a permutation over $\mathbb{F}_{2^{n}}$ then $F\left(x^{3}\right)$ is APN on $\mathbb{F}_{2^{n}}$.\\
Proof: Since the function $F^{\prime}(x)=F\left(x^{3}\right)$ is quadratic then it is APN if and only if, for any $a \in \mathbb{F}_{2^{n}}^{*}$, the only solutions of the linear homogeneous equation


\begin{equation*}
F^{\prime}(a x)+F^{\prime}(a x+a)+F^{\prime}(a)=0 \tag{1}
\end{equation*}


are 0,1 . We get the equation

$$
a^{3}\left(x+x^{2}\right)+L\left(a^{9}\left(x^{8}+x\right)\right)=0
$$

and denoting $y=x+x^{2}$, we obtain

$$
a^{3} y+L\left(a^{9}\left(y^{4}+y^{2}+y\right)\right)=0
$$

with $\operatorname{tr}_{n}(y)=0$. Hence, denoting $u=a^{3} y$ and

$$
A_{u}=\left\{L\left(u^{3}\left(y+\frac{1}{y}+\frac{1}{y^{2}}\right)\right): y \in \mathbb{F}_{2^{n}}^{*}, \operatorname{tr}_{n}(y)=0\right\},
$$

we see that if the condition


\begin{equation*}
u \notin A_{u} \text { for every } u \in \mathbb{F}_{2^{n}}^{*} \tag{2}
\end{equation*}


is satisfied then $F^{\prime}$ is APN. Note that $A_{u} \subseteq\left\{L\left(u^{3} z\right): z \in\right. \left.\mathbb{F}_{2^{n}}, \operatorname{tr}_{n}(z)=0\right\}$.

Since $F$ is a permutation then $F(u x)+F(u x+u) \neq 0$ for any $x$ and any nonzero $u$, that is, $u+L\left(u^{3}\left(x^{2}+x+1\right)\right) \neq 0$. Thus for any nonzero $u$


\begin{equation*}
u \notin\left\{L\left(u^{3} v\right): v \in \mathbb{F}_{2^{n}}, \operatorname{tr}_{n}(v)=\operatorname{tr}_{n}(1)\right\} . \tag{3}
\end{equation*}


For $n$ even $\operatorname{tr}_{n}(1)=0$ and

$$
u \notin A_{u} \subseteq\left\{L\left(u^{3} v\right): v \in \mathbb{F}_{2^{n}}, \operatorname{tr}_{n}(v)=\operatorname{tr}_{n}(1)\right\} .
$$

Hence, condition (2) is satisfied and $F^{\prime}$ is indeed APN.\\
Remark 1: Note that condition (2) has to be checked only for the elements $u$ in the image of $L$. For instance, in the case $L(x)=\operatorname{tr}_{n}(x)$, this condition has to be checked for $u=1$ only and it is obviously satisfied. This gives a simpler way of\\
proving that $x^{3}+\operatorname{tr}_{n}\left(x^{9}\right)$ is APN for every $n$ (it is also easy to show that $x+\operatorname{tr}_{n}\left(x^{3}\right)$ is bijective).

Thus, to construct APN functions using Proposition 1, one has to look for permutations of the form $x+L\left(x^{3}\right)$ for a linear function $L$. It happens that examples of such permutations have been already found [2], in the framework of APN functions but in a different context. Note, that if $x+L\left(x^{3}\right)$ is a permutation then $x+a^{-1} L\left(a^{3} x^{3}\right)$ is also a permutation for every $a \neq 0$. Up to this equivalence, there are three known classes of such permutations:

\begin{itemize}
  \item $x+\operatorname{tr}_{n}\left(x^{3}\right)$ for any even $n$.
  \item $x+\operatorname{tr}_{n}^{3}\left(x^{12}+x^{6}\right)$ for $n$ divisible by 6 .
  \item $x+\operatorname{tr}_{n}^{3}\left(x^{3}+x^{6}\right)$ for $n$ divisible by 6 .
\end{itemize}

We used Proposition 5 given below to search for more permutations of the desired type with the help of a computer. Using Proposition 5 a quite efficient backtracking search can be implemented. This approach allows to search for all such permutations up to $n=14$ within a few days. Surprisingly, it turns out that the three classes listed above already cover all the cases for $n<16, n$ even. We pose the following open problem.

Open Problem 1: Find more permutations of the form $x+ L\left(x^{3}\right)$ not listed above or prove that the list above is already complete.

Next, we apply Proposition 1 to the second and the third classes of permutations listed above to construct APN functions (the first class gives the known APN function $x^{3}+\operatorname{tr}_{n}\left(x^{9}\right)$ ).

Corollary 1: Let $n$ be a positive integer divisible by 6 . Then the following functions from $\mathbb{F}_{2^{n}}$ to itself are APN


\begin{align*}
& x^{3}+\operatorname{tr}_{n}^{3}\left(x^{18}+x^{36}\right),  \tag{4}\\
& x^{3}+\operatorname{tr}_{n}^{3}\left(x^{9}+x^{18}\right) . \tag{5}
\end{align*}


We extend now slightly Proposition 1. Let $n$ and $L, L^{\prime}$ be linear functions from $\mathbb{F}_{2^{n}}$ to itself. According to Remark 3 of [2], if the function $F(x)=L^{\prime}(x)+L\left(x^{3}\right)$ is a permutation over $\mathbb{F}_{2^{n}}$ then $L^{\prime}$ is a permutation. Hence, by Proposition 1, if $F$ is a permutation then $F\left(x^{3}\right)=L^{\prime}\left(x^{3}\right)+L\left(x^{9}\right)$ is APN. This fact is formulated in the following proposition.

Proposition 2: Let $n$ be even and $L$ and $L^{\prime}$ be linear functions from $\mathbb{F}_{2^{n}}$ to itself. If $F(x)=L^{\prime}(x)+L\left(x^{3}\right)$ is a permutation over $\mathbb{F}_{2^{n}}$ then $F\left(x^{3}\right)$ is APN on $\mathbb{F}_{2^{n}}$.

For instance, for $n$ even, let $l(x)=a x+b$ with $a \in \mathbb{F}_{2^{n}}^{*}$, $b \in \mathbb{F}_{2^{n}}$, and let $L$ be a linear function from $\mathbb{F}_{2^{n}}$ to itself. If $F(x)=x+L\left(x^{3}\right)$ is a permutation then

$$
\begin{aligned}
F(l(x)) & =l(x)+L\left((l(x))^{3}\right) \\
& =a x+b+L\left(a^{3} x^{3}+a^{2} b x^{2}+a b^{2} x+b^{3}\right) \\
& =a x+L\left(a^{2} b x^{2}+a b^{2} x\right)+L\left(a^{3} x^{3}\right)+b+L\left(b^{3}\right)
\end{aligned}
$$

is a permutation too, and we can apply Proposition 2. We deduce:

Corollary 2: Let $n$ be even, let $a \in \mathbb{F}_{2^{n}}^{*}, b \in \mathbb{F}_{2^{n}}$, and let $L$ be a linear function from $\mathbb{F}_{2^{n}}$ to itself. If $x+L\left(x^{3}\right)$ is a permutation over $\mathbb{F}_{2^{n}}$ then the function

$$
a x^{3}+L\left(a^{3} x^{9}+a^{2} b x^{6}+a b^{2} x^{3}\right)
$$

is APN on $\mathbb{F}_{2^{n}}$.\\
Surprisingly enough, this almost straightforward corollary of Proposition 1 gives new examples of APN functions for different choices of elements $a$ and $b$ as we see it below.

Corollary 3: Let $a \in \mathbb{F}_{2^{n}}^{*}$. For $n$ even the function


\begin{equation*}
x^{3}+a^{-1} \operatorname{tr}_{n}\left(a^{3} x^{9}\right) \tag{6}
\end{equation*}


and for $n$ divisible by 6 the functions


\begin{align*}
& x^{3}+a^{-1} \operatorname{tr}_{n}^{3}\left(a^{6} x^{18}+a^{12} x^{36}\right)  \tag{7}\\
& x^{3}+a^{-1} \operatorname{tr}_{n}^{3}\left(a^{3} x^{9}+a^{6} x^{18}\right) \tag{8}
\end{align*}


are APN over $\mathbb{F}_{2^{n}}$.\\
Let $\alpha$ be a primitive element of $\mathbb{F}_{2^{n}}^{*}$ and $a=\alpha^{3 i+j}$ for some $i$ and $0 \leq j \leq 2$. Then $x^{3}+a^{-1} \operatorname{tr}_{n}\left(a^{3} x^{9}\right)=a^{-1}\left(\alpha^{j} y^{3}+\right. \left.\operatorname{tr}_{n}\left(\alpha^{3 j} y^{9}\right)\right)$ where $y=\alpha^{i} x$. If $j=2$ then taking $y=z^{2}$ we get $x^{3}+a^{-1} \operatorname{tr}_{n}\left(a^{3} x^{9}\right)=a^{-1}\left(\alpha z^{3}+\operatorname{tr}_{n}\left(\alpha^{3} z^{9}\right)\right)^{2}$. Hence, (6) can give at most two different cases $a=1$ and $a=\alpha$ (since all other cases are EA-equivalent to these two as shown above). This is also true for functions (7) and (8), that is, each of them can give at most two different cases $a=1$ and $a=\alpha$. An important observation is that for $a=1$ and $a=\alpha$ those functions are not necessarily CCZ-equivalent. For example, consider $n=6,8$, then the APN function $x^{3}+\operatorname{tr}_{n}\left(x^{9}\right)$ is CCZ-inequivalent to $x^{3}+\alpha^{-1} \operatorname{tr}_{n}\left(\alpha^{3} x^{9}\right)$. Hence, function (6) defines a new class of APN functions which includes the class of functions $x^{3}+\operatorname{tr}_{n}\left(x^{9}\right)$ from [3].

\section*{B. The case $n$ odd}
For $n$ even the function $x+\operatorname{tr}_{n}\left(x^{3}\right)$ and, when $n$ is also divisible by 3 , the functions $x+\operatorname{tr}_{n}^{3}\left(x^{3}+x^{6}\right)$ and $x+\operatorname{tr}_{n}^{3}\left(x^{6}+x^{12}\right)$ are permutations while they are not permutations for $n$ odd. However, as we shall see further, the corresponding functions $x^{3}+\operatorname{tr}_{n}\left(x^{9}\right), x^{3}+\operatorname{tr}_{n}^{3}\left(x^{9}+x^{18}\right)$ and $x^{3}+\operatorname{tr}_{n}^{3}\left(x^{18}+x^{36}\right)$ are APN for both $n$ even and $n$ odd cases. Besides, Proposition 1 does not apply to the case $n$ odd. However, it can be adapted to this case. Assume first that $F(x)=x+L\left(x^{3}\right)$ is a permutation over the super-field $\mathbb{F}_{2^{2 n}}$, then $F\left(x^{3}\right)$ is APN over $\mathbb{F}_{2^{n}}$, since we know then that, for every nonzero element $u$ of $\mathbb{F}_{2^{n}}$, we have $u \notin\left\{L\left(u^{3} v\right), v \in\right. \left.\mathbb{F}_{2^{2 n}}, \operatorname{tr}_{2 n}(v)=0\right\}$ and that these $v$ include all the elements of $\mathbb{F}_{2^{n}}$ and a fortiori those of null trace $\operatorname{tr}_{n}$. In fact, thanks to the observations in the proof of Proposition 1, this condition can be weakened and generalized:

Proposition 3: Let $n$ be any positive integer and $K$ some field extension of $\mathbb{F}_{2^{n}}$. Let $L$ be an $\mathbb{F}_{2}$-linear mapping from $\mathbb{F}_{2^{n}}$ to $\mathbb{F}_{2^{n}}$ extended to an $\mathbb{F}_{2}$-linear mapping from $K$ to $K$. Let $E$ be a coset in $K$ of a vector space containing $L\left(\mathbb{F}_{2^{n}}\right)$. Assume that $F(x)=x+L\left(x^{3}\right)$ is injective on $E$ and that the set $\left\{x^{2}+x+1 ; x \in E\right\}$ contains the set of elements $y$ of $\mathbb{F}_{2^{n}}$ such that $\operatorname{tr}_{n}(y)=0$. Then $F\left(x^{3}\right)$ is APN over $\mathbb{F}_{2^{n}}$.

We can take $K=\mathbb{F}_{2^{2 n}}$ and $E=w+\mathbb{F}_{2^{n}}$ where $w \in \mathbb{F}_{4} \backslash \mathbb{F}_{2}$. If $F(x)=x+L\left(x^{3}\right)$ is injective on $E$ then $F\left(x^{3}\right)$ is APN on $\mathbb{F}_{2^{n}}$ since the set $\left\{x^{2}+x+1=x+w+(x+w)^{2} ; x \in\right. E\}$ contains the set of elements $y$ of $\mathbb{F}_{2^{n}}$ such that $\operatorname{tr}_{n}(y)=0$.

For instance, if $n$ is odd and divisible by 3 then the functions

$$
\begin{aligned}
& F_{1}(x)=x+\operatorname{tr}_{n}^{3}\left(x^{6}+x^{12}\right), \\
& F_{2}(x)=x+\operatorname{tr}_{n}^{3}\left(x^{3}+x^{6}\right)
\end{aligned}
$$

are injective over $w+\mathbb{F}_{2^{n}} \subset \mathbb{F}_{2^{2 n}}$, where $w \in \mathbb{F}_{4} \backslash \mathbb{F}_{2}$. Indeed, the function $F_{2}$ is injective over $w+\mathbb{F}_{2^{n}}$ if and only if the function $x+w+\operatorname{tr}_{n}^{3}\left((x+w)^{3}+(x+w)^{6}\right)$ is injective over $\mathbb{F}_{2^{n}}$. We have

$$
\begin{aligned}
& x+w+\operatorname{tr}_{n}^{3}\left((x+w)^{3}+(x+w)^{6}\right) \\
= & x+w+\operatorname{tr}_{n}^{3}\left(x^{3}+x^{6}\right)+\operatorname{tr}_{n}^{3}\left(w^{2} x+w^{2} x^{4}+w^{3}+w^{6}\right) \\
= & x+w+\operatorname{tr}_{n}^{3}\left(x^{3}+x^{6}\right)+\operatorname{tr}_{n}^{3}\left(w^{2}\left(x+x^{4}\right)\right)
\end{aligned}
$$

because $w^{4}=w$ an $w^{3}=w^{6}=1$. Assume that there exist $x, a \in \mathbb{F}_{2^{n}}, a \neq 0$ such that

$$
\begin{aligned}
& x+\operatorname{tr}_{n}^{3}\left(x^{3}+x^{6}\right)+\operatorname{tr}_{n}^{3}\left(w^{2}\left(x+x^{4}\right)\right) \\
= & \left.(x+a)+\operatorname{tr}_{n}^{3}\left((x+a)^{3}+(x+a)^{6}\right)+\operatorname{tr}_{n}^{3}\left(w^{2}\left((x+a)+(x+a)^{4}\right)\right)\right)
\end{aligned}
$$

Then we get\\
$a+\operatorname{tr}_{n}^{3}\left(a x^{2}+a^{2} x+a^{2} x^{4}+a^{4} x^{2}+a^{3}+a^{6}\right)+\operatorname{tr}_{n}^{3}\left(w^{2}\left(a+a^{4}\right)\right)=0$.

Note that since $n$ is odd and $a \in \mathbb{F}_{2^{n}}$ then $\left(\operatorname{tr}_{n}^{3}\left(w^{2}(a+\right.\right. \left.\left.\left.a^{4}\right)\right)\right)^{8}=\operatorname{tr}_{n}^{3}\left(w\left(a+a^{4}\right)^{8}\right)=\operatorname{tr}_{n}^{3}\left(w^{2^{n-3}}\left(a+a^{4}\right)\right)= \operatorname{tr}_{n}^{3}\left(w\left(a+a^{4}\right)\right)$. Then adding (9) to its 8-th power we obtain $a+a^{8}+\operatorname{tr}_{n}^{3}\left(\left(w^{2}+w\right)\left(a+a^{4}\right)\right)=0$ which is actually $a+a^{8}+\operatorname{tr}_{n}^{3}\left(a+a^{4}\right)=0$ since $w^{2}+w=1$. Thus, $a+a^{8} \in \mathbb{F}_{8}$ and $a+a^{8}=\left(a+a^{8}\right)^{8}=a^{8}+a^{64}$, which gives $a=a^{64}$. Hence, $a \in \mathbb{F}_{2^{6}} \cap \mathbb{F}_{2^{n}}^{*}$, and since $n$ is odd then $a \in \mathbb{F}_{2^{3}}^{*}$ and $a+a^{3}+a^{6}+\operatorname{tr}_{n}^{3}\left(a^{2} x+a^{2} x^{4}+\left(a+a^{4}\right) x^{2}\right)+\left(a+a^{4}\right) \operatorname{tr}_{n}^{3}\left(w^{2}\right)=0$.\\
Since $\operatorname{tr}_{n}^{3}\left(w^{2}\right)=w^{i}$, where $i$ equals either 1 or 2 , then the previous equality gives\\
$a+a^{3}+a^{6}+\operatorname{tr}_{n}^{3}\left(a^{2} x+a^{2} x^{4}+\left(a+a^{4}\right) x^{2}\right)=\left(a+a^{4}\right) w^{i}$.\\
The left side of the equality above is in $\mathbb{F}_{8}$, the right in $w^{i} \cdot \mathbb{F}_{8}$. Hence, $a+a^{4}=0$ and $a \in \mathbb{F}_{4} \cap \mathbb{F}_{2^{3}}^{*}$, and, therefore, $a=1$. Then\\
$0=a+a^{3}+a^{6}+\operatorname{tr}_{n}^{3}\left(a^{2} x+a^{2} x^{4}+\left(a+a^{4}\right) x^{2}\right)=1+\operatorname{tr}_{n}^{3}\left(x+x^{4}\right)$.\\
The equation $1+\operatorname{tr}_{n}^{3}\left(x+x^{4}\right)=0$ has no solutions in $\mathbb{F}_{2^{n}}$. Indeed, if


\begin{equation*}
\operatorname{tr}_{n}^{3}\left(x+x^{4}\right)=1 \tag{10}
\end{equation*}


then


\begin{equation*}
\operatorname{tr}_{n}^{3}\left(x+x^{2}\right)=\left(\operatorname{tr}_{n}^{3}\left(x+x^{4}\right)\right)^{2}=1, \tag{11}
\end{equation*}


and adding (10) and (11)\\
$0=\operatorname{tr}_{n}^{3}\left(x+x^{4}+x+x^{2}\right)=\operatorname{tr}_{n}^{3}\left(x^{2}+x^{4}\right)=\left(\operatorname{tr}_{n}^{3}\left(x+x^{2}\right)\right)^{2}=1$.\\
A contradiction.\\
Thus, the function $F_{2}$ is indeed injective over $w+\mathbb{F}_{2^{n}} \subset \mathbb{F}_{2^{2 n}}$, and the proof for the case of the function $F_{1}$ is similar. Therefore, by Proposition 3 the functions $F_{1}\left(x^{3}\right)$ and $F_{2}\left(x^{3}\right)$ are APN over $\mathbb{F}_{2^{n}}$ for $n$ odd as well. Now we can formulate the following result.

Corollary 4: Let $n$ be a positive integer and $a \in \mathbb{F}_{2^{n}}^{*}$. Then function (6) is APN over $\mathbb{F}_{2^{n}}$, and, if $n$ is divisible by 3 , functions (7) and (8) are APN over $\mathbb{F}_{2^{n}}$.

Unlike to the case $n$ even, for $n$ odd different choices of $a \in \mathbb{F}_{2^{n}}^{*}$ for function (6) (respectively, functions (7) and (8)) give functions EA-equivalent to the case $a=1$. Indeed, $\operatorname{gcd}(n, 3)=1$ for $n$ odd and, for example, $a x^{3}+\operatorname{tr}_{n}\left(a^{3} x^{9}\right)= y^{3}+\operatorname{tr}_{n}\left(y^{9}\right)$ with $y=a^{\frac{1}{3}} x$.

We checked with a computer (by testing the isomorphy of the associated binary codes, see [5]) that, for $n=9$, functions (6), (7) and (8) are pairwise CCZ-inequivalent, and, that (7) and (8) are CCZ-inequivalent to any function from previously known families of APN functions. Hence, functions (7) and (8) define two new families of APN functions.

\section*{III. Sufficient conditions for permutations $x+L\left(x^{3}\right)$}
To construct new APN functions thanks to Proposition 1, we need to find more linear mappings $L$ such that $x+L\left(x^{3}\right)$ is a permutation. There is a simple sufficient condition for that.

Proposition 4: Let $n$ be a positive integer and $L$ be a linear function from $\mathbb{F}_{2^{n}}$ to itself. If for every $u \in \mathbb{F}_{2^{n}}$ such that $L(u) \neq 0$, the condition

$$
\operatorname{tr}_{n}\left(\frac{u}{(L(u))^{3}}\right)=\left\{\begin{array}{l}
0 \text { if } n \text { is odd } \\
1 \text { if } n \text { is even }
\end{array}\right.
$$

is satisfied, then $F(x)=x+L\left(x^{3}\right)$ is a permutation.\\
Proof: Let $x, y \in \mathbb{F}_{2^{n}}$ be any elements such that


\begin{equation*}
y=x+L\left(x^{3}\right) \tag{12}
\end{equation*}


Then, denoting $L^{i}(x)=(L(x))^{i}, i=2,3$, we have

$$
x^{3}=\left(y+L\left(x^{3}\right)\right)^{3}=y^{3}+y^{2} L\left(x^{3}\right)+y L^{2}\left(x^{3}\right)+L^{3}\left(x^{3}\right) .
$$

This means that there exists $z \in \mathbb{F}_{2^{m}}$ such that:


\begin{equation*}
z=(y+L(z))^{3}=y^{3}+y^{2} L(z)+y L^{2}(z)+L^{3}(z) . \tag{13}
\end{equation*}


Then, $F$ is a permutation if and only if, for every $y$, Equation (13) has at most one solution $z$. This is clear if $n$ is odd since the function $x \rightarrow x^{3}$ being then bijective, Relation (13) is equivalent to (12) with $z=x^{3}$. If $n$ is even, then if for some $y$, there exist several $z$ satisfying (13), then for each of them, $x=y+L(z)$ has $z$ for cube and is therefore a solution of equation $y=x+L\left(x^{3}\right)$, which means that $F$ is not injective; conversely, if for some $y$ there exist several $x$ satisfying (12), these elements cannot have same cube because of (12) and there are several $z$ satisfying (13).\\
A sufficient condition for equation (13) to have at most one solution for each $y$ is that, for every $u \neq 0$, there does not exist $y$ and $z$ such that

$$
u=y^{2} L(u)+y L^{2}(u)+L(z) L^{2}(u)+L^{2}(z) L(u)+L^{3}(u)
$$

since if $z$ and $z+u$ are solutions of (13), then this equality is satisfied. Note that $L(u)$ cannot be null since $u \neq 0$. We deduce that if, for every $u \neq 0$ such that $L(u) \neq 0$, there do not exist $y$ and $z$ such that

$$
\frac{u}{L^{3}(u)}=\left(\frac{y}{L(u)}\right)^{2}+\frac{y}{L(u)}+\left(\frac{L(z)}{L(u)}\right)^{2}+\frac{L(z)}{L(u)}+1
$$

then $F$ is a permutation. Hence a sufficient condition for $F$ to be a permutation is $\operatorname{tr}_{n}\left(\frac{u}{L^{3}(u)}\right)=1$, if $n$ is even and $\operatorname{tr}_{n}\left(\frac{u}{L^{3}(u)}\right)=0$, if $n$ is odd.

We shall give now a necessary and sufficient condition for $x+L\left(x^{3}\right)$ being bijective which is in particular helpful for a computer search. This proposition is actually a special case of Proposition 6 of [2] for which we give a slightly different proof. We first need to recall some definitions and facts. The Hamming weight $w t(f)$ of a Boolean function $f: \mathbb{F}_{2^{n}} \rightarrow \mathbb{F}_{2}$ is the size of its support $\left\{x \in \mathbb{F}_{2^{n}}: f(x) \neq 0\right\}$. A Boolean function $f$ is called balanced if $w t(f)=2^{n-1}$. A Boolean function $f$ is called bent if

$$
\lambda_{f}(a):=\sum_{x \in \mathbb{F}_{2^{n}}}(-1)^{f(x)+\operatorname{tr}_{n}(a x)}= \pm 2^{\frac{n}{2}}
$$

for any $a \in \mathbb{F}_{2^{n}}$. Bent functions are never balanced since $f$ is balanced if and only if $\lambda_{f}(0)=0$.

Proposition 5: Let $n$ be an even integer and $L$ be a linear function from $\mathbb{F}_{2^{n}}$ to itself. The function $x+L\left(x^{3}\right)$ is a permutation of $\mathbb{F}_{2^{n}}$ if and only if, for every $b \in \mathbb{F}_{2^{n}}^{*}$, such that $L^{*}(b) \neq 0$ there exists an element $\gamma \in \mathbb{F}_{2^{n}}$ such that $L^{*}(b)=\gamma^{3}$ and $\operatorname{tr}_{n}^{2}\left(\gamma^{-1} b\right) \neq 0$, were $L^{*}$ denotes the adjoint linear mapping of $L$.

Proof: The mapping $F(x)=L\left(x^{3}\right)+x$ is a permutation if and only if, for every $b \neq 0$, the Boolean function $\operatorname{tr}_{n}(b F(x))$ is balanced. We calculate

$$
\begin{aligned}
\lambda_{F}(0, b) & :=\sum_{x}(-1)^{\operatorname{tr}_{n}\left(b\left(L\left(x^{3}\right)+x\right)\right)} \\
& =\sum_{x}(-1)^{\operatorname{tr}_{n}\left(L^{*}(b) x^{3}+b x\right)} \\
& =\sum_{x}(-1)^{\operatorname{tr}_{n}\left(\frac{L^{*}(b)}{b^{3}} x^{3}+x\right)}=\lambda_{x^{3}}\left(1, \frac{L^{*}(b)}{b^{3}}\right) .
\end{aligned}
$$

Now we know that, for every $\beta \in \mathbb{F}_{2^{n}}$, the mapping $\operatorname{tr}_{n}\left(\beta x^{3}\right)$ is bent if and only if $\beta$ is a non-cube (see [7]), and this implies that $\lambda_{x^{3}}(1, \beta)$ is not zero in this case. Thus a necessary condition for $L\left(x^{3}\right)+x$ to be a permutation is that $\frac{L^{*}(b)}{b^{3}}$ is a cube for every $b$. This is equivalent to $L^{*}(b)$ being a cube for every $b$. Denote by $\gamma$ any third root of $L^{*}(b)$, i.e. $\gamma^{3}=L^{*}(b)$. The next step is to actually compute the value $\lambda_{x^{3}}(1, \beta)$ in this case. To simplify notation we denote $L^{*}(b) / b^{3}$ by $\beta$. Then

$$
\begin{aligned}
\lambda_{x^{3}}(1, \beta)^{2}= & \sum_{x, y}(-1)^{\operatorname{tr}_{n}\left(\beta x^{3}+x+\beta y^{3}+y\right)} \\
= & \sum_{y}(-1)^{\operatorname{tr}_{n}\left(\beta y^{3}+y\right)} \sum_{x}(-1)^{\operatorname{tr}_{n}\left(\beta\left(y^{2} x+y x^{2}\right)\right)} \\
& (\text { by replacing } y \text { by } x+y) \\
= & \sum_{y}(-1)^{\operatorname{tr}_{n}\left(\beta y^{3}+y\right)} \sum_{x}(-1)^{\operatorname{tr}_{n}\left(\left(\beta^{2} y^{4}+\beta y\right) x^{2}\right)} .
\end{aligned}
$$

Now the sum over x is 0 if $\beta^{2} y^{4}+\beta y \neq 0$ and $2^{n}$ otherwise. We denote $B=\gamma^{-1} b \cdot \mathbb{F}_{4}$. Since $\left(\gamma b^{-1}\right)^{3}=\beta$ then $\beta^{2} y^{4}+$\\
$\beta y=0$ if and only if $y \in B$. Then, using $\operatorname{tr}_{n}(1)=0$, we get

$$
\begin{aligned}
\lambda_{x^{3}}(1, \beta)^{2} & =2^{n} \sum_{y \in B}(-1)^{\operatorname{tr}_{n}\left(\beta y^{3}+y\right)}=2^{n} \sum_{y \in B}(-1)^{\operatorname{tr}_{n}(y)} \\
& =2^{n} \sum_{y \in \mathbb{F}_{4}}(-1)^{\operatorname{tr}_{n}\left(\gamma^{-1} b y\right)} \\
& =2^{n} \sum_{y \in \mathbb{F}_{4}}(-1)^{\operatorname{tr}_{2}\left(\operatorname{tr}_{n}^{2}\left(\gamma^{-1} b\right) y\right)} \\
& =\left\{\begin{array}{cc}
2^{n+2} & \text { if } \quad \operatorname{tr}_{n}^{2}\left(\gamma^{-1} b\right)=0 \\
0 & \text { else }
\end{array}\right.
\end{aligned}
$$

This completes the proof.\\
A very similar characterization holds in the case $n$ odd.\\
Proposition 6: Let $n$ be an odd integer and $L$ be a linear function from $\mathbb{F}_{2^{n}}$ to itself. The function $x+L\left(x^{3}\right)$ is a permutation of $\mathbb{F}_{2^{n}}$ if and only if for all $b \in \mathbb{F}_{2^{n}}$ either $L^{*}(b)=0$ or $\operatorname{tr}_{n}\left(\gamma^{-1} b\right)=0$, were $L^{*}(b)=\gamma^{3}$ and $L^{*}$ denotes the adjoint linear mapping of $L$.

Proof: Again, the mapping $F(x)=L\left(x^{3}\right)+x$ is a permutation if and only if, for every $b \neq 0$, the Boolean function $\operatorname{tr}_{n}(b F(x))$ is balanced. A similar calculation as in the proof of Proposition 5 yields

$$
\begin{aligned}
\lambda_{F}(0, b) & =\sum_{x}(-1)^{\operatorname{tr}_{n}\left(b\left(L\left(x^{3}\right)+x\right)\right)} \\
& =\sum_{x}(-1)^{\operatorname{tr}_{n}\left(\left(L^{*}(b)^{1 / 3} x\right)^{3}+b x\right)} \\
& =\sum_{x}(-1)^{\operatorname{tr}_{n}\left(x^{3}+\frac{b}{L^{*}(b)^{1 / 3}} x\right)} \\
& =\lambda_{x^{3}}\left(\frac{b}{L^{*}(b)^{1 / 3}}, 1\right)
\end{aligned}
$$

Now in the case $n$ odd $\lambda_{x^{3}}(a, 1)=0$ if and only if $\operatorname{tr}_{n}(a)=0$ (see for example [4]) and thus $F$ is a permutation if and only if $\operatorname{tr}_{n}\left(\frac{b}{L^{*}(b)^{1 / 3}}\right)=0$ for all $b$ such that $L^{*}(b) \neq 0$. This completes the proof.

A little more can be said in the case where all coefficients of $L$ are in the subfield $\mathbb{F}_{2^{n / 2}}$, as it is the case for all the examples we know. For this we first state a technical lemma.

Lemma 1: Let $n$ be an even integer and $L$ a linear function from $\mathbb{F}_{2^{n}}$ to itself with $L(x)=\sum_{i=0}^{n-1} \lambda_{i} x^{2^{i}}$ where all $\lambda_{i} \in \mathbb{F}_{2^{n / 2}}$. Furthermore denote by $L^{\prime}$ the restriction of $L$ to $\mathbb{F}_{2^{n / 2}}$. It holds that $L^{\prime *}(x)=L^{*}(x)$ for all $x \in \mathbb{F}_{2^{n / 2}}$, i.e. the restriction of the adjoint operator of $L$ to $\mathbb{F}_{2^{n / 2}}$ is the adjoint operator of the restriction of $L$ to $\mathbb{F}_{2^{n / 2}}$.

Proof: First note that when all $\lambda_{i}$ are in the subfield, so are the coefficients of $L^{*}$. Now let $x, y \in \mathbb{F}_{2^{n / 2}}$ be arbitrary and choose $z \in \mathbb{F}_{2^{n}}$ such that $y=z^{2^{n / 2}}+z$. We compute

$$
\begin{aligned}
\operatorname{tr}_{n / 2}\left(x L^{\prime}(y)\right) & =\operatorname{tr}_{n / 2}(x L(y))=\operatorname{tr}_{n / 2}\left(x L\left(z+z^{2^{n / 2}}\right)\right) \\
& =\operatorname{tr}_{n / 2}\left(x L(z)+(x L(z))^{2^{n / 2}}\right) \\
& =\operatorname{tr}_{n}(x L(z))=\operatorname{tr}_{n}\left(L^{*}(x) z\right) \\
& =\operatorname{tr}_{n / 2}\left(L^{*}(x)\left(z+z^{2^{n / 2}}\right)\right) \\
& =\operatorname{tr}_{n / 2}\left(L^{*}(x) y\right)
\end{aligned}
$$

As $L^{\prime *}$ is the unique linear mapping fulfilling $\operatorname{tr}_{n / 2}\left(x L^{\prime}(y)\right)= \operatorname{tr}_{n / 2}\left(L^{\prime *}(x) y\right)$ the lemma follows.

Proposition 7: Let $n$ be an even integer and $L$ be a linear function from $\mathbb{F}_{2^{n}}$ to itself such that the function $x+L\left(x^{3}\right)$ is a permutation. If all coefficients of $L$ are in $\mathbb{F}_{2^{n / 2}}$ then $L$ fulfills $L\left(\mathbb{F}_{2^{n / 2}}\right)=\{0\}$ and $L^{2}=0$.

Proof: In this case $F\left(\mathbb{F}_{2^{n / 2}}\right) \subseteq \mathbb{F}_{2^{n / 2}}$ and therefore $F$ is also bijective on this subfield. We first prove that $L^{*}(b)=0$ for all $b \in \mathbb{F}_{2^{n / 2}}$ and therefore that $L(b)=0$ for all $b \in \mathbb{F}_{2^{n / 2}}$. We consider the two cases $n / 2$ even and odd separately. If $n / 2$ is even then applying Proposition 5 on $x+L\left(x^{3}\right)$ viewed as a permutation on $\mathbb{F}_{2^{n}}$ gives $L^{*}(b)=0$ or $L^{*}(b)=\gamma^{3}$ and $\operatorname{tr}_{n}^{2}\left(\gamma^{-1} b\right) \neq 0$. But for $b \in \mathbb{F}_{2^{n / 2}}$ the second condition is impossible as when $\gamma^{3}$ is in $\mathbb{F}_{2^{n / 2}}$ so is $\gamma$ and therefore $\operatorname{tr}_{n}^{2}\left(\gamma^{-1} b\right)=\operatorname{tr}_{n / 2}^{2}\left(\gamma^{-1} b \operatorname{tr}_{n}^{n / 2}(1)\right)=0$. Thus we see that $L^{*}(b)=0$ for all $b \in \mathbb{F}_{2^{n / 2}}$. In the case $n / 2$ odd we apply Proposition 6 on $x+L\left(x^{3}\right)$ viewed as a permutation on $\mathbb{F}_{2^{n / 2}}$ (and for $b \in \mathbb{F}_{2^{n / 2}}$ ) and conclude that either $L^{*}(b)=0$ or $L^{*}(b)=\gamma^{3}$ where $\operatorname{tr}_{n / 2}\left(\gamma^{-1} b\right)=0$. On the other hand applying Proposition 5 on $x+L\left(x^{3}\right)$ viewed as a permutation on $\mathbb{F}_{2^{n}}$ and Lemma 1 yields in the case where $L^{*}(b) \neq 0$ that $\operatorname{tr}_{n}^{2}\left(\gamma^{-1} b\right) \neq 0$, a contradiction as for an element $a \in \mathbb{F}_{2^{n / 2}}$ we have $\operatorname{tr}_{n / 2}(a)=\operatorname{tr}_{n}^{2}(a)$. Thus for all even $n$ it holds that $L^{*}(b)=0$ for all $b \in \mathbb{F}_{2^{n / 2}}$. This implies that $L(b)=0$ for all $b \in \mathbb{F}_{2^{n / 2}}$. Now, in both cases ( $n$ even and $n$ odd), we choose $x \in \mathbb{F}_{2^{n}}$ arbitrarily and compute $L(x)+L(x)^{2^{n / 2}}=L(x)+L\left(x^{2^{n / 2}}\right)=L\left(x+x^{2^{n / 2}}\right)=0$ which is equivalent to $L(x) \in \mathbb{F}_{2^{n / 2}}$. Hence, $L(L(x))=0$ for any $x \in \mathbb{F}_{2^{n}}$.

Finally, we recall a characterization of permutation of the form $F(x)=L\left(x^{3}\right)+L^{\prime}(x)$ given in [2].

Proposition 8 (Prop. 5 of [2]): Let $L$ and $L^{\prime}$ be linear functions from $\mathbb{F}_{2^{n}}$ to itself. The mapping $F(x)=L\left(x^{3}\right)+ L^{\prime}(x)$ is a permutation over $\mathbb{F}_{2^{n}}$ if and only if, for every $u \in \mathbb{F}_{2^{n}}^{*}$ and every $v$ such that $\operatorname{tr}_{n}(v)=\operatorname{tr}_{n}(1)$, the condition $L\left(u^{3} v\right) \neq L^{\prime}(u)$ holds.

\section*{References}
[1] E. Biham and A. Shamir. Differential Cryptanalysis of DES-like Cryptosystems. Journal of Cryptology, vol. 4, No.1, pp. 3-72, 1991.\\[0pt]
[2] L. Budaghyan, C. Carlet, A. Pott. New Classes of Almost Bent and Almost Perfect Nonlinear Functions. IEEE Trans. Inform. Theory, vol. 52, no. 3, pp. 1141-1152, March 2006.\\[0pt]
[3] L. Budaghyan, C. Carlet, G. Leander. Constructing new APN functions from known ones. Finite Fields and Their Applications, v. 15, issue 2, pp. 150-159, April 2009.\\[0pt]
[4] C. Carlet. Vectorial Boolean Functions for Cryptography. Chapter of the monography Boolean Methods and Models, Y. Crama and P. Hammer eds, Cambridge University Press, to appear soon. Preliminary version available at \href{http://www-rocq.inria.fr/codes/Claude.Carlet/pubs.html}{http://www-rocq.inria.fr/codes/Claude.Carlet/pubs.html}\\[0pt]
[5] C. Carlet, P. Charpin and V. Zinoviev. Codes, bent functions and permutations suitable for DES-like cryptosystems. Designs, Codes and Cryptography, 15(2), pp. 125-156, 1998.\\[0pt]
[6] F. Chabaud and S. Vaudenay. Links between differential and linear cryptanalysis, Advances in Cryptology -EUROCRYPT'94, LNCS, SpringerVerlag, New York, 950, pp. 356-365, 1995.\\[0pt]
[7] G. Leander. Monomial bent functions. IEEE Transactions on Information Theory, vol. 52, no. 2, pp. 738-743, 2006.\\[0pt]
[8] M. Matsui. Linear cryptanalysis method for DES cipher. Advances in Cryptology-EUROCRYPT'93, LNCS, Springer-Verlag, pp. 386-397, 1994.\\[0pt]
[9] K. Nyberg. Differentially uniform mappings for cryptography, Advances in Cryptography, EUROCRYPT'93, LNCS, Springer-Verlag, New York, 765, pp. 55-64, 1994.


\end{document}